%%%%%%%%%%%%%%%%%%%%%%%%%%%%%%%%%%%%%%%%%%%%%%%%%%%%%%%%%%%%%%%%%%%%%%%%%%%%%%%%
% PREÂMBULO EXTENSO
%%%%%%%%%%%%%%%%%%%%%%%%%%%%%%%%%%%%%%%%%%%%%%%%%%%%%%%%%%%%%%%%%%%%%%%%%%%%%%%%
\documentclass[12pt, a4paper]{article}

% --- Pacotes de Codificação, Fonte e Língua ---
\usepackage[utf8]{inputenc}
\usepackage[T1]{fontenc}
\usepackage{lmodern}
\usepackage[portuguese]{babel}

% --- Pacotes de Layout e Geometria ---
\usepackage[margin=1in]{geometry} % Margens de 1 polegada (aprox 2.54cm)
\usepackage[parskip=half]{parskip} % Espaço entre parágrafos, em vez de indentação
\usepackage{setspace} % Para controlar o espaçamento entre linhas
% \onehalfspacing % Descomente para espaçamento 1.5

% --- Pacotes de Matemática Avançada ---
\usepackage{amsmath} % Essencial para equações (align, gather, pmatrix)
\usepackage{amssymb} % Símbolos matemáticos (como \mathbb{R}, \approx)
\usepackage{amsthm}  % Para ambientes de Teorema/Definição

% --- Pacotes de Estilo e Utilidades ---
\usepackage{xcolor} % Para definir cores personalizadas
\usepackage{mdframed} % Para criar caixas de destaque
\usepackage{booktabs} % Para tabelas de qualidade profissional (ex: \toprule)
\usepackage{graphicx} % Para incluir imagens (se necessário)
\usepackage{float}    % Para melhor controle de flutuantes (figuras, tabelas)
\usepackage{hyperref} % Para links clicáveis (Sumário, referências)

% --- Configuração do Hyperref ---
\definecolor{linkblue}{RGB}{0,0,180}
\hypersetup{
    colorlinks=true,
    linkcolor=linkblue,
    filecolor=magenta,      
    urlcolor=cyan,
    pdftitle={Resumo Detalhado - Estimação por Intervalo},
    pdfauthor={Miqueias Teixeira},
    bookmarks=true,
    pdfpagemode=UseOutlines
}

% --- Definição de Novos Ambientes (Teoremas, Definições, etc.) ---
% \newtheorem{nome_ambiente}{Título que aparecerá}[contador_seção]
\newtheorem{teorema}{Teorema}[section]
\newtheorem{definicao}{Definição}[section]
\newtheorem{propriedade}{Propriedade}[section]
\newtheorem{corolario}{Corolário}[section]

% --- Ambientes de Estilo Próprio (para Exemplos, Interpretações) ---
\theoremstyle{definition} % Muda o estilo para ter fonte normal, não itálico
\newtheorem{exemplo}{Exemplo Numérico}[section]
\newtheorem{aviso}{Cuidado!}[section]

\theoremstyle{remark} % Estilo de "nota"
\newtheorem{interpretacao}{Interpretação}[section]
\newtheorem{analogia}{Analogia}[section]

% --- Estilo da Caixa de "Receita" ---
\mdfdefinestyle{caixa_receita}{
    frametitle={A Receita de Bolo:},
    frametitlebackgroundcolor=gray!20,
    frametitlerule=true,
    backgroundcolor=gray!5,
    innertopmargin=10pt,
    innerbottommargin=10pt,
    linewidth=1pt,
    linecolor=black!70,
    skipabove=10pt,
    skipbelow=10pt,
    nobreak=true % Tenta não quebrar a caixa entre páginas
}

\mdfdefinestyle{caixa_bayes}{
    frametitle={A Receita Bayesiana:},
    frametitlebackgroundcolor=blue!10,
    frametitlerule=true,
    backgroundcolor=blue!5,
    innertopmargin=10pt,
    innerbottommargin=10pt,
    linewidth=1pt,
    linecolor=blue!70,
    skipabove=10pt,
    skipbelow=10pt,
    nobreak=true
}

% --- Título e Autor ---
\title{Guia de Estudo Detalhado \\ \Large Capítulo 5: Estimação por Intervalo}
\author{Compilado para Miqueias Teixeira \\ \small (Baseado no Capítulo 5 de seu material)}
\date{\today}

%%%%%%%%%%%%%%%%%%%%%%%%%%%%%%%%%%%%%%%%%%%%%%%%%%%%%%%%%%%%%%%%%%%%%%%%%%%%%%%%
% DOCUMENTO
%%%%%%%%%%%%%%%%%%%%%%%%%%%%%%%%%%%%%%%%%%%%%%%%%%%%%%%%%%%%%%%%%%%%%%%%%%%%%%%%
\begin{document}

\maketitle
\tableofcontents
\newpage

% -------------------------------------------------------------------
% SEÇÃO 1: A GRANDE IDEIA
% -------------------------------------------------------------------
\section{A Grande Ideia: A Incerteza da Estimativa}

\subsection{O Problema da Estimativa Pontual}
Quando queremos saber um parâmetro $\theta$ da população (como a média $\mu$ ou a variância $\sigma^2$), não podemos medir todos. Então, pegamos uma amostra $X_1, \dots, X_n$.

A partir da amostra, calculamos uma \textbf{estimativa pontual}, $\hat{\theta}$. Por exemplo, usamos $\bar{X}$ para estimar $\mu$.

\begin{aviso}
    Qual a chance de $\bar{X}$ ser \textit{exatamente} igual a $\mu$?
    $$ P(\bar{X} = \mu) = 0 $$
    (Para variáveis contínuas). Nossa estimativa pontual está \textit{quase certamente errada}. Ela pode estar perto, mas não é exata. O trabalho de um estatístico é \textbf{quantificar a incerteza} desse "perto".
\end{aviso}

\subsection{A Solução: O Intervalo de Confiança}
Em vez de dar um único número (um "chute"), damos uma faixa de valores plausíveis.

\begin{analogia}{A Rede de Pesca}
    \begin{itemize}
        \item \textbf{Estimativa Pontual:} É como tentar pegar um peixe ($\mu$) com um arpão ($\hat{\theta}$). Você tem que acertar \textit{exatamente}. É muito provável que você erre.
        \item \textbf{Estimativa por Intervalo:} É como usar uma \textbf{rede de pesca} ($[L(X), U(X)]$). Você joga a rede na água (calcula o intervalo). Você não sabe onde o peixe está, mas você tem uma certa "confiança" de que sua rede o capturou.
    \end{itemize}
\end{analogia}

\begin{definicao}{Intervalo de Confiança (IC)}
    Um IC para $\theta$ com \textbf{nível de confiança} $\gamma = 1-\alpha$ (ex: 95\%, onde $\alpha=0.05$) é um intervalo \textbf{aleatório} $[L(X), U(X)]$, onde $L(X)$ e $U(X)$ são estatísticas (calculadas da amostra), tal que:
    $$ P(L(X) \le \theta \le U(X)) = \gamma $$
\end{definicao}

\begin{interpretacao}{A Verdade sobre os 95\%}
    Essa é a parte mais mal compreendida da estatística.
    \begin{itemize}
        \item \textbf{Errado:} "Há 95\% de chance do verdadeiro $\mu$ estar dentro do intervalo [10, 20]."
        \item \textbf{Correto:} "Eu tenho 95\% de confiança neste \textit{método}. Se eu repetisse meu experimento (coletar a amostra e construir o intervalo) 100 vezes, eu esperaria que cerca de 95 dos meus intervalos \textit{capturassem} o verdadeiro valor de $\mu$."
    \end{itemize}
    Uma vez que você calcula o intervalo, $[10, 20]$, ele não é mais aleatório. O parâmetro $\mu$ (que é fixo) ou \textbf{está} lá dentro ou \textbf{não está}. O 95\% é a nossa confiança no \textit{processo} que gerou o intervalo, não em um intervalo específico.
\end{interpretacao}

\newpage
% -------------------------------------------------------------------
% SEÇÃO 2: A FERRAMENTA MÁGICA (PIVOTAL)
% -------------------------------------------------------------------
\section{A Receita de Bolo: O Método da Quantidade Pivotal (Seção 5.2)}

Como encontramos $L(X)$ e $U(X)$? Precisamos de uma "fórmula mágica" que conecte nossa amostra ao parâmetro.

\begin{definicao}{Quantidade Pivotal (Pivô)}
    Uma \textbf{Quantidade Pivotal}, $Q(X, \theta)$, é uma função da amostra $X = (X_1, \dots, X_n)$ e do parâmetro $\theta$ que satisfaz duas condições \textit{cruciais}:
    \begin{enumerate}
        \item A fórmula \textbf{contém} o parâmetro $\theta$ que queremos estimar.
        \item A \textbf{distribuição de probabilidade} de $Q(X, \theta)$ é \textbf{completamente conhecida} e \textbf{NÃO depende} de $\theta$ (ou de qualquer outro parâmetro desconhecido).
    \end{enumerate}
    A distribuição do pivô é "universal" (ex: $N(0,1)$, $\chi^2_{df}$, $t_{df}$). Ela só pode depender de constantes conhecidas, como o tamanho da amostra $n$.
\end{definicao}

\begin{mdframed}[frametitle={Receita de 3 Passos: Construindo um IC com um Pivô}, style=caixa_receita]
    Para construir um IC de $\gamma = 1-\alpha$ para $\theta$:

    \textbf{Passo 1: Encontrar o Pivô} \\
    Encontre uma quantidade pivotal $Q(X, \theta)$ e sua distribuição (ex: $N(0,1)$, $t_{n-1}$, etc.).

    \textbf{Passo 2: Encontrar os Quantis (Cortes)} \\
    Vá na tabela da distribuição do Pivô e encontre dois valores (quantis), $q_1$ e $q_2$, tais que a área entre eles seja $\gamma$:
    $$ P(q_1 \le Q(X, \theta) \le q_2) = \gamma $$
    \begin{itemize}
        \item \textbf{Simétrico:} $P(Q < q_1) = \alpha/2$ e $P(Q > q_2) = \alpha/2$.
        \item \textbf{Assimétrico (ex: $\chi^2$):} Os valores $q_1$ e $q_2$ não serão simétricos (ex: $q_1 \neq -q_2$).
    \end{itemize}
    
    \textbf{Passo 3: Isolar $\theta$ (A Mágica da Álgebra)} \\
    Pegue a inequação do Passo 2 e resolva-a algebricamente para isolar $\theta$ no meio.
    $$ q_1 \le Q(X, \theta) \le q_2 \quad \xrightarrow{\text{álgebra}} \quad L(X) \le \theta \le U(X) $$
    O resultado $[L(X), U(X)]$ é o seu intervalo de confiança.
\end{mdframed}

% -------------------------------------------------------------------
% SEÇÃO 3: A CAIXA DE FERRAMENTAS (DISTRIBUIÇÕES)
% -------------------------------------------------------------------
\section{A Caixa de Ferramentas Amostral (\texorpdfstring{$Z$, $T$, $\chi^2$, $F$}{Z, T, chi^2, F})}
Antes de aplicar a receita, vamos conhecer nossas ferramentas (os pivôs). Todas elas nascem de amostras da $N(0,1)$.

\begin{propriedade}{As 4 Distribuições da Inferência}
    Seja $Z \sim N(0,1)$ e $Z_1, \dots, Z_k \sim N(0,1)$ independentes.
    \begin{enumerate}
        \item \textbf{Normal (Z):} O "Pai" de todos. Usado quando $\sigma^2$ é conhecida.
            $$ Z = \frac{\bar{X} - \mu}{\sigma / \sqrt{n}} \sim N(0,1) $$
        \item \textbf{Qui-Quadrado ($\chi^2$):} A "Mãe da Variância". É a soma de normais padrão ao quadrado.
            $$ U = \sum_{i=1}^k Z_i^2 \sim \chi^2_k \quad (\text{k graus de liberdade}) $$
        \item \textbf{t-Student (T):} O "Z humilde". É uma Normal (Z) dividida pela raiz de uma $\chi^2$. Usado quando \textit{não} sabemos $\sigma^2$.
            $$ T = \frac{Z}{\sqrt{U/k}} \sim t_k \quad (\text{k graus de liberdade}) $$
        \item \textbf{F-Snedecor (F):} A "Juíza". É a razão de duas $\chi^2$. Usada para \textit{comparar} duas variâncias.
            $$ F = \frac{U_1 / k_1}{U_2 / k_2} \sim F_{k_1, k_2} \quad (\text{$k_1, k_2$ graus de liberdade}) $$
    \end{enumerate}
\end{propriedade}

O \textbf{Teorema 5.1} do seu livro é a ponte que nos permite usar essas ferramentas.

\begin{teorema}[Teorema 5.1 - A Ponte da População Normal]
    Seja $X_1, \dots, X_n$ uma amostra i.i.d. de uma população $N(\mu, \sigma^2)$. Então:
    \begin{enumerate}
        \item[(i)] $\bar{X}$ e $S^2$ são variáveis aleatórias \textbf{independentes}.
        \item[(ii)] $Q = \frac{(n-1)S^2}{\sigma^2} \sim \chi^2_{n-1}$.
    \end{enumerate}
\end{teorema}
\textbf{Por que isso é genial?}
\begin{itemize}
    \item A Parte (ii) nos dá \textit{exatamente} um Pivô para $\sigma^2$. É uma $\chi^2$ (Ferramenta 2).
    \item A Parte (i) e (ii) juntas nos permitem \textit{provar} que $\frac{\bar{X} - \mu}{S / \sqrt{n}}$ é uma $t_{n-1}$ (Ferramenta 3).
\end{itemize}

\newpage
% -------------------------------------------------------------------
% SEÇÃO 4: ICs PARA UMA AMOSTRA
% -------------------------------------------------------------------
\section{Aplicando a Receita: ICs para Uma Amostra (Seção 5.3)}

\subsection{\texorpdfstring{IC para $\mu$ (com $\sigma^2$ conhecida) - O ``Caso Z''}{IC para mu (com sigma^2 conhecida) - O ``Caso Z''}}
\begin{itemize}
    \item \textbf{Contexto:} Raro, mas fundamental. Ex: Você está em uma fábrica, a máquina \textit{por design} tem um desvio padrão $\sigma=2$mm, mas você acha que a média $\mu$ descalibrou.
    \item \textbf{Passo 1 (Pivô):} $Q(X, \mu) = Z = \frac{\bar{X} - \mu}{\sigma / \sqrt{n}} \sim N(0,1)$.
    \item \textbf{Passo 2 (Quantis):} Para $\gamma = 1-\alpha$, queremos $q_1 = -z_{\alpha/2}$ e $q_2 = +z_{\alpha/2}$.
        $$ P(-z_{\alpha/2} \le Z \le z_{\alpha/2}) = 1-\alpha $$
    \item \textbf{Passo 3 (Álgebra):}
    \begin{align*}
        -z_{\alpha/2} &\le \frac{\bar{X} - \mu}{\sigma / \sqrt{n}} \le z_{\alpha/2} \\
        -z_{\alpha/2} \left(\frac{\sigma}{\sqrt{n}}\right) &\le \bar{X} - \mu \le z_{\alpha/2} \left(\frac{\sigma}{\sqrt{n}}\right) \\
        -\bar{X} - z_{\alpha/2} \left(\frac{\sigma}{\sqrt{n}}\right) &\le - \mu \le -\bar{X} + z_{\alpha/2} \left(\frac{\sigma}{\sqrt{n}}\right) \\
        \bar{X} + z_{\alpha/2} \left(\frac{\sigma}{\sqrt{n}}\right) &\ge \mu \ge \bar{X} - z_{\alpha/2} \left(\frac{\sigma}{\sqrt{n}}\right)
    \end{align*}
    \item \textbf{Resultado (IC para $\mu$, $\sigma^2$ conhecida):}
        $$ IC(\mu, 1-\alpha) = \left[ \bar{X} - z_{\alpha/2} \frac{\sigma}{\sqrt{n}}, \quad \bar{X} + z_{\alpha/2} \frac{\sigma}{\sqrt{n}} \right] $$
\end{itemize}

\begin{exemplo}{Caso Z}
    Uma máquina de encher pacotes de café tem $\sigma=10$g (conhecido). Uma amostra de $n=25$ pacotes tem uma média amostral $\bar{X}=505$g. Construa um IC de 95\% para a média $\mu$.

    \textbf{1. Coletar Dados:}
    \begin{itemize}
        \item $\bar{X} = 505$g
        \item $\sigma = 10$g (conhecido)
        \item $n = 25$
        \item $\gamma = 0.95 \implies \alpha = 0.05 \implies \alpha/2 = 0.025$
    \end{itemize}
    \textbf{2. Achar Quantil:}
    \begin{itemize}
        \item Da tabela Normal, $z_{\alpha/2} = z_{0.025} = 1.96$.
    \end{itemize}
    \textbf{3. Calcular Margem de Erro (ME):}
    \begin{itemize}
        \item $ME = z_{\alpha/2} \frac{\sigma}{\sqrt{n}} = 1.96 \cdot \frac{10}{\sqrt{25}} = 1.96 \cdot \frac{10}{5} = 1.96 \cdot 2 = 3.92$
    \end{itemize}
    \textbf{4. Construir IC:}
    \begin{itemize}
        \item $IC(\mu, 95\%) = [\bar{X} - ME, \bar{X} + ME]$
        \item $IC(\mu, 95\%) = [505 - 3.92, \quad 505 + 3.92] = [501.08, \quad 508.92]$
    \end{itemize}
    \begin{interpretacao}
        Temos 95\% de confiança de que o \textit{método} que usamos gerou um intervalo que captura a verdadeira média $\mu$ de enchimento da máquina. Baseado nesta amostra, valores plausíveis para $\mu$ estão entre 501.08g e 508.92g. (Note que se a máquina devesse encher com 500g, o valor 500 está \textit{fora} do intervalo, sugerindo que a máquina está descalibrada).
    \end{interpretacao}
\end{exemplo}

\subsection{\texorpdfstring{IC para $\mu$ (com $\sigma^2$ desconhecida) - O ``Caso T''}{IC para mu (com sigma^2 desconhecida) - O ``Caso T''}}
\begin{itemize}
    \item \textbf{Contexto:} O caso mais comum. Você não sabe $\mu$ e também não sabe $\sigma^2$.
    \item \textbf{Passo 1 (Pivô):} Substituímos $\sigma$ por sua estimativa $S$. A distribuição muda de Z para T.
        $$ Q(X, \mu) = T = \frac{\bar{X} - \mu}{S / \sqrt{n}} \sim t_{n-1} $$
    \item \textbf{Passo 2 (Quantis):} Da tabela t-Student com $gl = n-1$.
        $$ P(-t_{n-1, \alpha/2} \le T \le t_{n-1, \alpha/2}) = 1-\alpha $$
    \item \textbf{Passo 3 (Álgebra):} Idêntica à do Caso Z.
    \item \textbf{Resultado (IC para $\mu$, $\sigma^2$ desconhecida):}
        $$ IC(\mu, 1-\alpha) = \left[ \bar{X} - t_{n-1, \alpha/2} \frac{S}{\sqrt{n}}, \quad \bar{X} + t_{n-1, \alpha/2} \frac{S}{\sqrt{n}} \right] $$
\end{itemize}

\begin{exemplo}{Caso T}
    Queremos saber o tempo médio $\mu$ de reação a um fármaco. Coletamos $n=16$ pacientes. A amostra deu $\bar{X}=2.5$s e $S=0.8$s. Construa um IC de 95\% para $\mu$.

    \textbf{1. Coletar Dados:}
    \begin{itemize}
        \item $\bar{X} = 2.5$s
        \item $S = 0.8$s (desconhecido, estimado)
        \item $n = 16$
        \item $\gamma = 0.95 \implies \alpha/2 = 0.025$
        \item Graus de Liberdade: $gl = n-1 = 15$
    \end{itemize}
    \textbf{2. Achar Quantil:}
    \begin{itemize}
        \item Da tabela t-Student, $t_{gl, \alpha/2} = t_{15, 0.025} = 2.131$.
        \item (Note que $t=2.131$ é $>$ $z=1.96$. A t-Student é mais "conservadora" e dá intervalos mais largos para compensar a incerteza de $S$).
    \end{itemize}
    \textbf{3. Calcular Margem de Erro (ME):}
    \begin{itemize}
        \item $ME = t_{gl, \alpha/2} \frac{S}{\sqrt{n}} = 2.131 \cdot \frac{0.8}{\sqrt{16}} = 2.131 \cdot \frac{0.8}{4} = 2.131 \cdot 0.2 = 0.4262$
    \end{itemize}
    \textbf{4. Construir IC:}
    \begin{itemize}
        \item $IC(\mu, 95\%) = [\bar{X} - ME, \bar{X} + ME]$
        \item $IC(\mu, 95\%) = [2.5 - 0.4262, \quad 2.5 + 0.4262] = [2.0738, \quad 2.9262]$
    \end{itemize}
    \begin{interpretacao}
        Temos 95\% de confiança que o verdadeiro tempo médio de reação $\mu$ está entre 2.07s e 2.93s.
    \end{interpretacao}
\end{exemplo}

\subsection{\texorpdfstring{IC para $\sigma^2$ - O ``Caso Qui-Quadrado''}{IC para sigma^2 - O ``Caso Qui-Quadrado''}}
\begin{itemize}
    \item \textbf{Contexto:} Queremos estimar a \textit{variabilidade} do processo, não a média.
    \item \textbf{Passo 1 (Pivô):} Do Teorema 5.1(ii):
        $$ Q(X, \sigma^2) = \frac{(n-1)S^2}{\sigma^2} \sim \chi^2_{n-1} $$
    \item \textbf{Passo 2 (Quantis):} \textbf{CUIDADO!} A $\chi^2$ não é simétrica.
        $$ P(q_1 \le Q \le q_2) = 1-\alpha $$
        \begin{itemize}
            \item $q_1 = \chi^2_{n-1, 1-\alpha/2}$ (Quantil da cauda \textit{esquerda})
            \item $q_2 = \chi^2_{n-1, \alpha/2}$ (Quantil da cauda \textit{direita})
        \end{itemize}
    \item \textbf{Passo 3 (Álgebra):}
    \begin{align*}
        q_1 &\le \frac{(n-1)S^2}{\sigma^2} \le q_2 \\
        \frac{1}{q_2} &\le \frac{\sigma^2}{(n-1)S^2} \le \frac{1}{q_1} \quad \text{(Invertemos tudo! Sinais trocam)} \\
        \frac{(n-1)S^2}{q_2} &\le \sigma^2 \le \frac{(n-1)S^2}{q_1}
    \end{align*}
    \item \textbf{Resultado (IC para $\sigma^2$):}
        $$ IC(\sigma^2, 1-\alpha) = \left[ \frac{(n-1)S^2}{\chi^2_{n-1, \alpha/2}}, \quad \frac{(n-1)S^2}{\chi^2_{n-1, 1-\alpha/2}} \right] $$
    \begin{aviso}
        Note a inversão! O quantil $q_2$ (direito, maior) vai para o limite \textbf{inferior}. O quantil $q_1$ (esquerdo, menor) vai para o limite \textbf{superior}.
    \end{aviso}
\end{itemize}

\begin{exemplo}{Caso Qui-Quadrado}
    Usando os dados do fármaco ($n=16$, $S=0.8$s), construa um IC de 95\% para a \textit{variância} $\sigma^2$ do tempo de reação.

    \textbf{1. Coletar Dados:}
    \begin{itemize}
        \item $S = 0.8 \implies S^2 = 0.64$
        \item $n = 16$
        \item $\gamma = 0.95 \implies \alpha = 0.05 \implies \alpha/2 = 0.025$ e $1-\alpha/2 = 0.975$
        \item Graus de Liberdade: $gl = n-1 = 15$
    \end{itemize}
    \textbf{2. Achar Quantis (da tabela $\chi^2$):}
    \begin{itemize}
        \item Cauda Esquerda: $q_1 = \chi^2_{15, 0.975} = 6.262$
        \item Cauda Direita: $q_2 = \chi^2_{15, 0.025} = 27.488$
    \end{itemize}
    \textbf{3. Calcular Limites (usando a fórmula invertida):}
    \begin{itemize}
        \item Numerador: $(n-1)S^2 = 15 \cdot (0.8)^2 = 15 \cdot 0.64 = 9.6$
        \item Limite Inferior: $L(X) = \frac{(n-1)S^2}{q_2} = \frac{9.6}{27.488} = 0.349$
        \item Limite Superior: $U(X) = \frac{(n-1)S^2}{q_1} = \frac{9.6}{6.262} = 1.533$
    \end{itemize}
    \textbf{4. Construir IC:}
    \begin{itemize}
        \item $IC(\sigma^2, 95\%) = [0.349, \quad 1.533]$
    \end{itemize}
    \begin{interpretacao}
        Temos 95\% de confiança que a verdadeira variância $\sigma^2$ do tempo de reação está entre 0.349 e 1.533.
        \textit{(Se quiséssemos um IC para o desvio padrão $\sigma$, bastaria tirar a raiz quadrada dos limites: $[\sqrt{0.349}, \sqrt{1.533}] \approx [0.591, 1.238]$).}
    \end{interpretacao}
\end{exemplo}

\newpage
% -------------------------------------------------------------------
% SEÇÃO 5: ICs PARA DUAS AMOSTRAS
% -------------------------------------------------------------------
\section{Aplicando a Receita: ICs para Duas Amostras (Seção 5.3.2)}

\subsection{\texorpdfstring{IC para $\mu_X - \mu_Y$ (Assumindo $\sigma_X^2 = \sigma_Y^2 = \sigma^2$)}{IC para mu_X - mu_Y (Assumindo sigma_X^2 = sigma_Y^2 = sigma^2)}}
\begin{itemize}
    \item \textbf{Contexto:} Queremos comparar a média de dois grupos (ex: Droga A vs. Placebo B). Assumimos que as variâncias dos grupos são \textit{iguais}.
    \item \textbf{Estratégia:} Como $\sigma^2$ é a mesma, usamos as duas amostras para estimá-la. Criamos a \textbf{variância agrupada (pooled)}:
        $$ S_p^2 = \frac{(n_X-1)S_X^2 + (n_Y-1)S_Y^2}{n_X + n_Y - 2} $$
    É uma média ponderada das variâncias amostrais.
    \item \textbf{Passo 1 (Pivô):} Generalização do Caso T.
        $$ Q = T = \frac{(\bar{X} - \bar{Y}) - (\mu_X - \mu_Y)}{S_p \sqrt{\frac{1}{n_X} + \frac{1}{n_Y}}} \sim t_{n_X+n_Y-2} $$
    \item \textbf{Passo 2 (Quantis):} Da tabela t com $gl = n_X + n_Y - 2$.
    \item \textbf{Passo 3 (Álgebra):} Isolar $(\mu_X - \mu_Y)$.
    \item \textbf{Resultado (IC para $\mu_X - \mu_Y$):}
        $$ IC(\mu_X - \mu_Y, 1-\alpha) = \left[ (\bar{X} - \bar{Y}) \pm t_{gl, \alpha/2} \cdot S_p \sqrt{\frac{1}{n_X} + \frac{1}{n_Y}} \right] $$
\end{itemize}

\begin{exemplo}{Caso T (Duas Amostras)}
    Grupo X (Droga): $n_X=10, \bar{X}=5.5, S_X=1.2$ \\
    Grupo Y (Placebo): $n_Y=12, \bar{Y}=4.8, S_Y=1.0$ \\
    Construa um IC de 90\% para $\mu_X - \mu_Y$ (assumindo $\sigma_X^2 = \sigma_Y^2$).

    \textbf{1. Coletar Dados:}
    \begin{itemize}
        \item $\bar{X} - \bar{Y} = 5.5 - 4.8 = 0.7$
        \item $\gamma = 0.90 \implies \alpha = 0.10 \implies \alpha/2 = 0.05$
        \item $gl = n_X + n_Y - 2 = 10 + 12 - 2 = 20$
    \end{itemize}
    \textbf{2. Calcular $S_p^2$ e $S_p$:}
    \begin{itemize}
        \item $S_p^2 = \frac{(10-1)(1.2)^2 + (12-1)(1.0)^2}{20} = \frac{9 \cdot 1.44 + 11 \cdot 1}{20}$
        \item $S_p^2 = \frac{12.96 + 11}{20} = \frac{23.96}{20} = 1.198$
        \item $S_p = \sqrt{1.198} \approx 1.0945$
    \end{itemize}
    \textbf{3. Achar Quantil:}
    \begin{itemize}
        \item $t_{gl, \alpha/2} = t_{20, 0.05} = 1.725$
    \end{itemize}
    \textbf{4. Calcular Margem de Erro (ME):}
    \begin{itemize}
        \item $ME = t_{gl, \alpha/2} \cdot S_p \sqrt{\frac{1}{n_X} + \frac{1}{n_Y}} = 1.725 \cdot 1.0945 \cdot \sqrt{\frac{1}{10} + \frac{1}{12}}$
        \item $ME = 1.888 \cdot \sqrt{0.1 + 0.0833} = 1.888 \cdot \sqrt{0.1833} = 1.888 \cdot 0.428 = 0.808$
    \end{itemize}
    \textbf{5. Construir IC:}
    \begin{itemize}
        \item $IC(\mu_X - \mu_Y, 90\%) = [ (\bar{X} - \bar{Y}) \pm ME ]$
        \item $IC = [0.7 - 0.808, \quad 0.7 + 0.808] = [-0.108, \quad 1.508]$
    \end{itemize}
    \begin{interpretacao}
        Temos 90\% de confiança de que a verdadeira diferença $\mu_X - \mu_Y$ está entre -0.108 e 1.508. Como o valor \textbf{zero está} dentro do intervalo, \textbf{não podemos concluir} (com 90\% de confiança) que há uma diferença entre a droga e o placebo.
    \end{interpretacao}
\end{exemplo}

\subsection{\texorpdfstring{IC para $\sigma_X^2 / \sigma_Y^2$ - O ``Caso F''}{IC para sigma_X^2 / sigma_Y^2 - O ``Caso F''}}
\begin{itemize}
    \item \textbf{Contexto:} Queremos comparar a \textit{variabilidade} de dois grupos.
    \item \textbf{Passo 1 (Pivô):} A "Juíza" F, que é a razão de duas $\chi^2$ (divididas por seus $gl$).
        $$ Q = F = \frac{S_X^2 / \sigma_X^2}{S_Y^2 / \sigma_Y^2} = \left(\frac{S_X^2}{S_Y^2}\right) \left(\frac{\sigma_Y^2}{\sigma_X^2}\right) \sim F_{n_X-1, n_Y-1} $$
    \item \textbf{Passo 2 (Quantis):} Assimétrico. $gl_1 = n_X-1, gl_2 = n_Y-1$.
        \begin{itemize}
            \item $q_1 = F_{gl_1, gl_2, 1-\alpha/2}$ (Quantil esquerdo)
            \item $q_2 = F_{gl_1, gl_2, \alpha/2}$ (Quantil direito)
        \end{itemize}
    \item \textbf{Passo 3 (Álgebra):} Isolar $\sigma_X^2 / \sigma_Y^2$.
    \begin{align*}
        q_1 &\le \left(\frac{S_X^2}{S_Y^2}\right) \left(\frac{\sigma_Y^2}{\sigma_X^2}\right) \le q_2 \\
        q_1 \cdot \left(\frac{S_Y^2}{S_X^2}\right) &\le \frac{\sigma_Y^2}{\sigma_X^2} \le q_2 \cdot \left(\frac{S_Y^2}{S_X^2}\right) \\
        \frac{1}{q_2 \cdot (S_Y^2/S_X^2)} &\ge \frac{\sigma_X^2}{\sigma_Y^2} \ge \frac{1}{q_1 \cdot (S_Y^2/S_X^2)}
    \end{align*}
    \item \textbf{Resultado (IC para $\sigma_X^2 / \sigma_Y^2$):}
        $$ IC\left(\frac{\sigma_X^2}{\sigma_Y^2}, 1-\alpha\right) = \left[ \frac{1}{q_2} \left(\frac{S_X^2}{S_Y^2}\right), \quad \frac{1}{q_1} \left(\frac{S_X^2}{S_Y^2}\right) \right] $$
\end{itemize}

\begin{exemplo}{Caso F}
    Usando os dados do Exemplo 4.2 ($n_X=10, S_X=1.2$; $n_Y=12, S_Y=1.0$), construa um IC de 95\% para $\sigma_X^2 / \sigma_Y^2$.

    \textbf{1. Coletar Dados:}
    \begin{itemize}
        \item $S_X^2 = 1.44$, $S_Y^2 = 1.0 \implies S_X^2 / S_Y^2 = 1.44$
        \item $gl_1 = n_X - 1 = 9$
        \item $gl_2 = n_Y - 1 = 11$
        \item $\gamma = 0.95 \implies \alpha/2 = 0.025$ e $1-\alpha/2 = 0.975$
    \end{itemize}
    \textbf{2. Achar Quantis (da tabela F):}
    \begin{itemize}
        \item Cauda Direita: $q_2 = F_{9, 11, 0.025} = 3.59$ (aprox.)
        \item Cauda Esquerda: $q_1 = F_{9, 11, 0.975}$. (Truque: $F_{a,b, 1-p} = 1 / F_{b,a, p}$)
        \item $q_1 = 1 / F_{11, 9, 0.025} = 1 / 3.96 \approx 0.252$
    \end{itemize}
    \textbf{3. Calcular Limites:}
    \begin{itemize}
        \item Limite Inferior: $L = \frac{1}{q_2} \left(\frac{S_X^2}{S_Y^2}\right) = \frac{1}{3.59} \cdot 1.44 = 0.401$
        \item Limite Superior: $U = \frac{1}{q_1} \left(\frac{S_X^2}{S_Y^2}\right) = \frac{1}{0.252} \cdot 1.44 = 5.714$
    \end{itemize}
    \textbf{4. Construir IC:}
    \begin{itemize}
        \item $IC(\sigma_X^2 / \sigma_Y^2, 95\%) = [0.401, \quad 5.714]$
    \end{itemize}
    \begin{interpretacao}
        Temos 95\% de confiança de que a verdadeira razão das variâncias está entre 0.401 e 5.714. Como o valor \textbf{um está} dentro do intervalo, \textbf{não podemos rejeitar} a hipótese de que as variâncias são iguais ($\sigma_X^2 = \sigma_Y^2$). Isso valida a premissa que usamos no Exemplo 4.2.
    \end{interpretacao}
\end{exemplo}

\newpage
% -------------------------------------------------------------------
% SEÇÃO 6: ICs APROXIMADOS (EMV)
% -------------------------------------------------------------------
\section{Intervalos de Confiança Aproximados (Seção 5.4)}

\subsection{Quando a População Não é Normal?}
Se $n$ é pequeno e a população \textit{não} é normal, estamos em apuros (precisamos de métodos não-paramétricos ou bootstrap).
Mas se $n$ é \textbf{grande} (ex: $n > 30$), podemos invocar o \textbf{Teorema Central do Limite (TCL)} e a teoria assintótica do \textbf{Estimador de Máxima Verossimilhança (EMV)}.

\begin{teorema}{Normalidade Assintótica do EMV}
    Seja $\hat{\theta}$ o EMV de $\theta$. Para $n$ grande:
    $$ \hat{\theta} \approx N\left(\theta, \frac{1}{I(\theta)}\right) $$
    Onde $I(\theta)$ é a \textbf{Informação de Fisher}. Na prática, usamos a \textbf{Informação de Fisher Observada}, $I(\hat{\theta})$, que é a Informação de Fisher avaliada no EMV.
\end{teorema}

\begin{itemize}
    \item \textbf{Passo 1 (Pivô Aproximado):} Padronizando o teorema acima.
        $$ Q = Z = \frac{\hat{\theta} - \theta}{\sqrt{1 / I(\hat{\theta})}} \approx N(0,1) $$
    \item \textbf{Passo 2 e 3:} Idênticos ao "Caso Z".
    \item \textbf{Resultado (IC Aproximado para $\theta$):}
        $$ IC(\theta, 1-\alpha) = \left[ \hat{\theta} \pm z_{\alpha/2} \frac{1}{\sqrt{I(\hat{\theta})}} \right] $$
\end{itemize}

\begin{exemplo}{IC Aproximado para Exponencial}
    Seja $X_i \sim \text{Exp}(\theta)$, onde $f(x) = \frac{1}{\theta}e^{-x/\theta}$. (Aqui, $\theta = E(X)$).
    Encontramos $\hat{\theta}_{EMV} = \bar{X}$.
    A Informação de Fisher é $I(\theta) = \frac{n}{\theta^2}$.
    A Informação Observada é $I(\hat{\theta}) = \frac{n}{\hat{\theta}^2} = \frac{n}{\bar{X}^2}$.
    O Pivô é: $Z = \frac{\bar{X} - \theta}{\sqrt{1 / (n/\bar{X}^2)}} = \frac{\bar{X} - \theta}{\bar{X} / \sqrt{n}} \approx N(0,1)$.

    \textbf{Dados:} Coletamos $n=50$ amostras e $\bar{X}=10.5$. IC de 95\%.
    
    \textbf{1. Coletar Dados:}
    \begin{itemize}
        \item $\hat{\theta} = \bar{X} = 10.5$
        \item $n = 50$
        \item $z_{\alpha/2} = 1.96$
    \end{itemize}
    \textbf{2. Calcular Erro Padrão (denominador):}
    \begin{itemize}
        \item $\text{ep}(\hat{\theta}) = \frac{1}{\sqrt{I(\hat{\theta})}} = \frac{1}{\sqrt{n/\bar{X}^2}} = \frac{\bar{X}}{\sqrt{n}}$
        \item $\text{ep}(\hat{\theta}) = \frac{10.5}{\sqrt{50}} = \frac{10.5}{7.071} = 1.485$
    \end{itemize}
    \textbf{3. Calcular IC:}
    \begin{itemize}
        \item $IC = [ \hat{\theta} \pm z_{\alpha/2} \cdot \text{ep}(\hat{\theta}) ]$
        \item $IC = [ 10.5 \pm 1.96 \cdot 1.485 ] = [ 10.5 \pm 2.91 ] = [7.59, \quad 13.41]$
    \end{itemize}
    \begin{interpretacao}
        Temos 95\% de confiança que a verdadeira média $\theta$ da distribuição Exponencial está entre 7.59 e 13.41. (Note que este intervalo \textit{não} é o mesmo que o intervalo $\chi^2$ exato do Exercício 5.6, mas é uma boa aproximação por $n$ ser grande).
    \end{interpretacao}
\end{exemplo}

\newpage
% -------------------------------------------------------------------
% SEÇÃO 7: INTERVALOS BAYESIANOS
% -------------------------------------------------------------------
\section{Uma Filosofia Diferente: Intervalos Bayesianos (Seção 5.5)}

\subsection{A Mudança de Paradigma: $\theta$ é Aleatório}
Na estatística Frequencista (Clássica), o parâmetro $\theta$ é um valor \textbf{fixo} e desconhecido. O intervalo $[L(X), U(X)]$ é \textbf{aleatório}.

Na estatística \textbf{Bayesiana}, o parâmetro $\theta$ \textbf{é uma variável aleatória}. Nós temos "crenças" sobre $\theta$. Os dados $X$ (que agora são fixos, pois foram observados) são usados para \textit{atualizar} nossa crença.

\begin{itemize}
    \item \textbf{Crença a Priori:} O que eu acho de $\theta$ \textit{antes} de ver os dados? $\pi(\theta)$.
    \item \textbf{Crença a Posteriori:} O que eu acho de $\theta$ \textit{depois} de ver os dados? $\pi(\theta | X)$.
    \item \textbf{Intervalo de Credibilidade:} Um intervalo que contém $\gamma$ (ex: 95\%) da massa de probabilidade da distribuição \textit{a posteriori}.
\end{itemize}

\begin{interpretacao}{Bayesiana vs. Frequencista}
    \begin{itemize}
        \item \textbf{Frequencista (IC de 95\%):} "Eu tenho 95\% de confiança no \textit{método} que gerou este intervalo."
        \item \textbf{Bayesiana (ICred de 95\%):} "Dados os dados que observei e minha crença inicial, há 95\% de \textit{probabilidade} do parâmetro $\theta$ estar dentro deste intervalo."
    \end{itemize}
    A interpretação Bayesiana é muito mais intuitiva, mas requer uma premissa inicial (a priori).
\end{interpretacao}

\begin{mdframed}[frametitle={A Receita Bayesiana: Construindo um Intervalo de Credibilidade}, style=caixa_bayes]
    \textbf{Passo 1: Definir a Priori, $\pi(\theta)$} \\
    Qual é a nossa crença sobre $\theta$ \textit{antes} de ver qualquer dado? (Ex: $\theta \sim \text{Gama}(a, b)$).

    \textbf{Passo 2: Definir a Verossimilhança (Likelihood), $L(\theta | X)$} \\
    Esta é a $f(X | \theta)$, a função de densidade dos dados, vista como uma função de $\theta$. (Ex: $L(\theta|X) \propto e^{-n\theta} \theta^{\sum x_i}$ para a Poisson).

    \textbf{Passo 3: Calcular a Posteriori, $\pi(\theta | X)$} \\
    Usamos o Teorema de Bayes para atualizar nossa crença:
    $$ \pi(\theta | X) \propto L(\theta | X) \cdot \pi(\theta) $$
    (Multiplicamos a Verossimilhança pela Priori).

    \textbf{Passo 4: Encontrar o Intervalo} \\
    Encontramos $q_1$ e $q_2$ que "cortam" a distribuição a posteriori:
    $$ \int_{q_1}^{q_2} \pi(\theta | X) d\theta = \gamma $$
    Para um intervalo simétrico (caudas iguais), $q_1$ é o quantil $\alpha/2$ e $q_2$ é o quantil $1-\alpha/2$ da \textit{distribuição a posteriori}.
\end{mdframed}

\begin{exemplo}{Bayesiano (do Exercício 5.11)}
    Dados $X_i \sim \text{Poisson}(\theta)$, $n=10$, $S = \sum X_i = 18$.
    Priori: $\pi(\theta) = e^{-\theta}, \theta > 0$. (Que é $\text{Gama}(\alpha=1, \beta=1)$).
    ICred de 95\%.

    \textbf{Passo 1 (Priori):}
    \begin{itemize}
        \item $\pi(\theta) \sim \text{Gama}(\alpha_{priori}=1, \beta_{priori}=1)$.
        \item Kernel: $\pi(\theta) \propto \theta^{1-1} e^{-\theta/1} = e^{-\theta}$.
    \end{itemize}
    \textbf{Passo 2 (Verossimilhança):}
    \begin{itemize}
        \item $L(\theta|X) = \prod \frac{e^{-\theta}\theta^{X_i}}{X_i!} \propto e^{-n\theta} \theta^{\sum X_i}$
        \item $L(\theta|X) \propto e^{-10\theta} \theta^{18}$.
    \end{itemize}
    \textbf{Passo 3 (Posteriori):}
    \begin{itemize}
        \item $\pi(\theta | X) \propto L(\theta | X) \cdot \pi(\theta)$
        \item $\pi(\theta | X) \propto (e^{-10\theta} \theta^{18}) \cdot (e^{-\theta})$
        \item $\pi(\theta | X) \propto e^{-11\theta} \theta^{18}$
    \end{itemize}
    \textbf{4. Reconhecer a Posteriori:}
    \begin{itemize}
        \item O kernel $e^{-11\theta} \theta^{18}$ é de uma $\text{Gama}(\alpha, \beta)$
        \item $\theta^{\alpha-1} e^{-\theta/\beta} \implies \alpha-1 = 18 \implies \alpha = 19$
        \item $1/\beta = 11 \implies \beta = 1/11$
        \item \textbf{A Posteriori é $\theta|X \sim \text{Gama}(19, 1/11)$}.
    \end{itemize}
    \textbf{5. Construir Intervalo:}
    \begin{itemize}
        \item Queremos $q_1$ e $q_2$ da $\text{Gama}(19, 1/11)$ que deixam $\alpha/2=0.025$ em cada cauda.
        \item $q_1 = Q_{\text{Gama}}(0.025 \mid \alpha=19, \beta=1/11) \approx 0.945$
        \item $q_2 = Q_{\text{Gama}}(0.975 \mid \alpha=19, \beta=1/11) \approx 2.701$
        \item (Estes valores vêm de software, ex: R `qgamma(c(0.025, 0.975), shape=19, scale=1/11)`)
        \item $ICred(\theta, 95\%) = [0.945, \quad 2.701]$
    \end{itemize}
    \begin{interpretacao}
        Dados os 18 eventos observados em 10 amostras, e nossa crença inicial de $\pi(\theta)=e^{-\theta}$, temos 95\% de probabilidade de que a verdadeira taxa $\theta$ esteja entre 0.945 e 2.701.
    \end{interpretacao}
\end{exemplo}

% -------------------------------------------------------------------
% SEÇÃO 8: COMPARAÇÃO E ARMADILHAS
% -------------------------------------------------------------------
\section{Resumo: Frequencista vs. Bayesiano e Armadilhas Comuns}

\subsection{Tabela de Comparação}

\begin{table}[H]
    \centering
    \caption{Comparação das Filosofias de Intervalo}
    \begin{tabular}{@{}lll@{}}
        \toprule
        Característica & Frequencista (Clássico) & Bayesiano \\ \midrule
        \textbf{O que é $\theta$?} & Fixo, desconhecido & Aleatório, com distribuição \\
        \textbf{O que é o IC?} & \textbf{Aleatório} & \textbf{Fixo} (após os dados) \\
        \textbf{Intervalo} & Intervalo de \textbf{Confiança} & Intervalo de \textbf{Credibilidade} \\
        \textbf{Interpretação} & "Confiança no \textit{método}" & "Probabilidade de $\theta$ \textit{estar no} intervalo" \\
        \textbf{Requer} & Um Pivô & Uma Priori $\pi(\theta)$ \\
        \textbf{Resultado} & $P(L(X) \le \theta \le U(X)) = \gamma$ & $P(q_1 \le \theta \le q_2 | X) = \gamma$ \\
        \bottomrule
    \end{tabular}
\end{table}

\subsection{Armadilhas Comuns (O que NÃO fazer)}
\begin{aviso}{Armadilha 1: A Interpretação}
    \textbf{NÃO DIGA:} "Há 95\% de chance do $\mu$ estar no IC de [10, 20]."
    O $\mu$ é fixo. Ele ou está (100\% de chance) ou não está (0\% de chance). O 95\% é sobre a \textit{taxa de sucesso do seu método} a longo prazo.
\end{aviso}

\begin{aviso}{Armadilha 2: O Nível de Confiança}
    Um IC de 99\% é \textbf{mais largo} (menos preciso) que um IC de 90\%.
    \begin{itemize}
        \item Para ter \textbf{mais confiança} (99\%), sua "rede de pesca" precisa ser \textbf{maior} para garantir que você pegue o peixe.
        \item Para ter \textbf{menos confiança} (90\%), você pode usar uma rede menor.
    \end{itemize}
\end{aviso}

\begin{aviso}{Armadilha 3: O Tamanho da Amostra $n$}
    A largura do intervalo é (geralmente) proporcional a $1/\sqrt{n}$.
    Para \textbf{diminuir} a largura do intervalo (ter mais precisão) pela \textbf{metade}, você precisa \textbf{quadruplicar} o tamanho da sua amostra ($n \to 4n$). A precisão custa caro!
\end{aviso}


\end{document}